\section{Introduction}

\subsection{Background}

%\noindent 写这一段只做三件事：现象是什么、危害在哪、为什么这一类方案适用。
% 不要在这里展开模型，也不要提前给结果——摘要已经给过了。
% 长度：4–6 句，一段到底，不分小点。
Large vehicles experience lateral vibration during operation ...
(symptom) ... It accelerates fatigue of the mounting structure and degrades
(positioning / comfort) ... Among the active schemes, centrifugal actuators
are suitable here because (why this class fits this problem) ...

\subsection{Work}

% 逐问复述"这道题要交什么"——一小节一问。
% 好处：评委读完 1.2 就知道三问各要交哪些曲线、哪些指标，
% 后面读到结果时不用回头找题目要求。
\subsubsection{Problem 1}
Under the operating conditions of Scenario 1 ... establish the lateral model,
report the displacement and acceleration curves for $0$--$10$\,s with the actuator
inactive and active, and give the vibration index for both states.

\subsubsection{Problem 2}
Under the multi-source disturbance of Scenario 2 ... design the angular displacement
schedule that minimises the index, report both states' curves and the schedule itself,
and verify the angular velocity / acceleration limits.

\subsubsection{Problem 3}
With the disturbance unmeasurable and only the lateral acceleration observable ...
design a schedule that minimises the index from the measurable signals alone, and
report the index for both states.

% ┈┈┈┈┈┈┈┈┈┈┈┈┈┈┈┈┈┈┈┈┈┈┈┈┈┈┈┈┈┈┈┈┈┈┈
% 【Section writing prompt】
% ════════════════════════════════════════════
% Position: the first section. 1.1 background is one paragraph, 1.2 repeats the tasks.
% 1.2 is not a summary of the paper: it restates what each question asks for. Write it in
%           the problem's own vocabulary, keep one subsubsection per question.
% Do not:   put method names in 1.1 unless the method is the subject; put results in 1.2.
% Citation: 1–2 references at most in 1.1 (the context claim), none in 1.2.
% ════════════════════════════════════════════
